\section{Primitive sufficient conditions}
\label{app:primitive}
\renewcommand{\theequation}{C.\arabic{equation}}
\setcounter{equation}{0}

Conditions (E1)--(E3) are stated on the projected quantiles. This
appendix gives primitive sufficient conditions --- conditions on
$\gamma$, on the full conditional survival function, and on the fibre
kernels $K_j(u,\cdot)$ --- that imply (E1), (E2) and the tail-only
condition (E3$^\uparrow$) of Appendix~\ref{app:aux}, with explicit
rates, and verifies them for the simulation models of
Section~\ref{sec:simulation}.

All constants and envelopes below are uniform in the triangular array,
in $j\le p_n$, and on the enlarged interval
$I_{\varepsilon/2}=[\varepsilon/2,1-\varepsilon/2]$; scores continue to
be integrated over $\I$. For $x=\iota_j(u,v)$ put
\begin{equation}\label{eq:deficit}
 \xi_j(u)=\esssup_{v\sim K_j(u,\cdot)}\gamma\{\iota_j(u,v)\},
 \qquad
 D_{j,u}(v)=\frac1{\gamma\{\iota_j(u,v)\}}-\frac1{\xi_j(u)}\ge0.
\end{equation}
The reciprocal-index deficit $D_{j,u}$, rather than
$\xi_j(u)-\gamma$, is what makes the Laplace representation below exact.

\subsection{Why (C1)--(C2) do not suffice}

\begin{proposition}[Necessary repairs]\label{prop:app-necessary}
The assumptions (C1)--(C2), even together with (S), imply none of (E1),
(E2), (E3).
\end{proposition}

\begin{proof}
(C1) constrains only the upper tail and does not exclude atoms, so it
does not imply the continuity of the projected conditional distribution
required by (E1).

For (E2), a level bound on a Hall remainder is insufficient. Let
$\bar F(e^t)=\exp\{-t+r(t)\}$, where $r$ is a sum of disjoint smooth
negative bumps centred at $t_m\to\infty$, the $m$th of height
$a_m=e^{-\beta t_m}$, with a decreasing side of width $o(a_m)$ and a
recovery side of slope at most $1/2$. Then $-1+r'(t)<0$, so $\bar F$ is
a valid survival function and $|r(t)|\le e^{-\beta t}$, yet
$r'(t_m)\to-\infty$. With $T(x)=\log Q(e^x)$, quantile inversion gives
\[
 \frac{\mathrm d}{\mathrm dx}\log\{Q(e^x)e^{-x}\}
 =\frac{r'\{T(x)\}}{1-r'\{T(x)\}},
\]
whose absolute value stays away from zero along a subsequence, so the
normalized increment in (E2) need not vanish.

Finally, make all conditional survival functions identical above a fixed
level but let the conditional median below it vary discontinuously with
the conditioning coordinate. This changes none of (C1), (C2), (S), and
none of the tail quantities, yet the printed (E3) fails because its
range $w\in[\tau_n,1-\tau_n]$ contains $w=1/2$. This last obstruction is
an artefact of the statement, not of the method: by
Lemma~\ref{lem:app-tailfreeze} the proofs use only $w\le3\alpha_n$, so
the correct repair is to weaken (E3) to (E3$^\uparrow$) rather than to
impose regularity on the body of the distribution.
\end{proof}

\subsection{The primitive conditions}

\begin{itemize}
\item[(P1)] \emph{Geometry and atomlessness.} For constants
 $0<\underline\gamma\le\Gamma<\infty$ and $L_\gamma<\infty$,
 \[
  \underline\gamma\le\gamma(x)\le\Gamma,
  \qquad
  |\gamma(x)-\gamma(x')|\le L_\gamma\|x-x'\|_\infty .
 \]
 Full fibre support (S) holds, and every conditional law $Y\mid U=x$ is
 atomless.
\item[(P2)] \emph{Differentiable uniform tail approximation.} For
 $t\ge t_0$,
 \[
  \Pp(Y>e^t\mid U=x)=c(x)e^{-t/\gamma(x)}\{1+r(t,x)\},
  \qquad 0<c_-\le c(x)\le c_+<\infty,
 \]
 with $\sup_x|r(t,x)|\le R_0(t)$ and
 $\sup_x|\partial_tr(t,x)|\le R_1(t)$, where $R_0(t)\to0$ and
 $\overline R_1(T)=\sup_{t\ge T}R_1(t)\to0$; eventually
 $R_0(t)\le1/2$. The maps $r(\cdot,x)$ are locally absolutely
 continuous.
\item[(P3)] \emph{Uniform doubling of the fibre-deficit mass.} With
 \[
  H_{j,u}(z)=\int\mathbf 1_{\{D_{j,u}(v)\le z\}}\,
  c\{\iota_j(u,v)\}\,K_j(u,\mathrm dv),
 \]
 there are $z_0,h_0>0$ and $C_D<\infty$ with
 \[
  H_{j,u}(z_0)\ge h_0,
  \qquad
  H_{j,u}(2z)\le C_DH_{j,u}(z)
  \quad(0<z\le z_0/2)
 \]
 uniformly. This allows an atom at $D=0$, includes the degenerate case
 $D\equiv0$ of an inactive coordinate, and rules out only fibres on
 which the probability of coming near the maximal index becomes
 arbitrarily thinner under a fixed rescaling.
\item[(P4)] \emph{Horizontal coupling of the primitive tail.} Whenever
 $u\in\I$, $v\in I_{\varepsilon/2}$ and $|u-v|\le\varepsilon/2$, the
 kernels $K_j(u,\cdot)$ and $K_j(v,\cdot)$ admit a coupling
 $X_u=\iota_j(u,V_u)$, $X_v=\iota_j(v,V_v)$ such that, almost surely and
 for every $t\ge t_0$,
 \[
  \left|\frac1{\gamma(X_v)}-\frac1{\gamma(X_u)}\right|\le L_D|v-u|,
 \]
 \[
  |\log c(X_v)-\log c(X_u)|
  +|\log\{1+r(t,X_v)\}-\log\{1+r(t,X_u)\}|
  \le L_H|v-u|+\zeta_n,
 \]
 with $\zeta_n\ge0$ and $\zeta_n=o(\Delta_{\min})$. This concerns only
 the primitive conditional tail of (P2), never the body.
\end{itemize}

A sufficient construction for (P4) is a coupling with
$\|X_v-X_u\|_\infty\le C_K|v-u|$ together with Lipschitz $1/\gamma$,
$\log c$ and $\log(1+r)$ on the coupled states; (P4) is weaker because it
records only the three quantities actually used.

\subsection{The transfer theorem}

\begin{lemma}[Tilted deficit bound]\label{lem:app-tilted}
Under (P3), for all sufficiently large $t$ the quantity
\[
 B_{j,u}(t)=\int e^{-tD_{j,u}(v)}c\{\iota_j(u,v)\}K_j(u,\mathrm dv)
\]
satisfies $B_{j,u}(t/2)\le C_BB_{j,u}(t)$ for a uniform constant $C_B$,
and consequently
\begin{equation}\label{eq:tiltedmoment}
 \frac{\int D_{j,u}(v)e^{-tD_{j,u}(v)}c\{\iota_j(u,v)\}K_j(u,\mathrm dv)}
      {B_{j,u}(t)}
 \;\le\;\frac{2C_B}{e\,t}
\end{equation}
uniformly in $n,j,u$.
\end{lemma}

\begin{proof}
Suppress $(j,u)$. For $t\ge z_0^{-1}$, restricting the integral to
$\{D\le1/t\}$ gives $B(t)\ge e^{-1}H(1/t)$. Split $B(t/2)$ over
$[0,1/t]$ and the dyadic bands $(2^m/t,2^{m+1}/t]$. Before the bands
reach $z_0$ their total contribution is at most
\[
 H(1/t)\Bigl[1+\sum_{m\ge0}e^{-2^{m-1}}C_D^{m+1}\Bigr],
\]
a convergent series. Iterating the doubling bound backwards from $z_0$
gives $H(1/t)\ge ct^{-\log_2C_D}$ for a uniform $c>0$, so the remaining
part, which starts no lower than $z_0/2$ and is at most
$c_+e^{-tz_0/4}$, is also bounded by a constant times $H(1/t)$.
Combining the two displays proves the doubling property of $B$. For
$z\ge0$, $ze^{-tz}\le(2/e)t^{-1}e^{-tz/2}$; integrating and using the
doubling property gives \eqref{eq:tiltedmoment}.
\end{proof}

\begin{theorem}[Primitive sufficient conditions]\label{thm:app-transfer}
Assume (P1)--(P4), let $\alpha_n\to0$, and suppose eventually
$C_0h_n\le\varepsilon/2$ and $\tau_n=(p_nn)^{-2}\le3\alpha_n<1$. Put
$x_\alpha=\log\{(3\alpha_n)^{-1}\}$ and $c_\gamma=\underline\gamma/2$.
Then, after a common deterministic onset,
\begin{align}
 \omega_\xi(r)&\le L_\gamma r, \label{eq:transfer-E1}\\
 b_n(3\alpha_n)&\le\frac{C_E}{x_\alpha}
 +4\Gamma^2\overline R_1(c_\gamma x_\alpha), \label{eq:transfer-E2}\\
 \omega_n^{\uparrow}(C_0h_n,\tau_n,\alpha_n)
 &\le2\Gamma C_0h_n\bigl[L_H+L_D\Gamma\{2\log(p_nn)+C_Q\}\bigr]
 +2\Gamma\zeta_n, \label{eq:transfer-E3}
\end{align}
where $C_E<\infty$ depends only on $\underline\gamma$, $\Gamma$, $C_B$
and the onset constants, and $C_Q=\max\{0,\log(2c_+)\}$. Consequently
(E1), (E2) and (E3$^\uparrow$) hold as soon as
\begin{equation}\label{eq:transfer-rates}
 n^{-1}=o(\Delta_{\min}),
 \quad
 x_\alpha^{-1}+\overline R_1(c_\gamma x_\alpha)=o(\Delta_{\min}),
 \quad
 h_n\log(p_nn)+\zeta_n=o(\Delta_{\min}).
\end{equation}
For a fixed gap and a Hall derivative envelope of order $O(t^{-1})$ or
smaller, the substantive requirement from (E2) is simply
$\log(1/\alpha_n)\to\infty$.
\end{theorem}

\begin{proof}
\emph{Step 1: (E1).} Under (S), $\xi_j(u)=\max_v\gamma\{\iota_j(u,v)\}$;
comparing the two maxima at the same $v$ in both directions gives
\eqref{eq:transfer-E1}, hence $\omega_\xi(1/n)\le L_\gamma/n$. A mixture
of atomless laws is atomless, because
$\Pp(Y=y\mid U_j=u)=\int\Pp\{Y=y\mid U=\iota_j(u,v)\}K_j(u,\mathrm dv)=0$;
this gives the continuity and the conditional quantile representation
required by (E1).

\emph{Step 2: exact projected factorization.} All derivative identities
below hold almost everywhere, and local absolute continuity justifies
their integrated forms. Let
\[
 \mathcal A_{j,u}(t)=\int e^{-tD_{j,u}(v)}c\{\iota_j(u,v)\}
 \{1+r(t,\iota_j(u,v))\}K_j(u,\mathrm dv),
\]
so that $\bar F_j(e^t\mid u)=e^{-t/\xi_j(u)}\mathcal A_{j,u}(t)$
exactly. With $M_{j,u}(t)=\int D_{j,u}e^{-tD_{j,u}}c\,K_j(u,\mathrm dv)$,
(P2), Lemma~\ref{lem:app-tilted} and $R_0\le1/2$ give
\begin{equation}\label{eq:app-g}
 \left|\frac{\mathcal A_{j,u}'(t)}{\mathcal A_{j,u}(t)}\right|
 \le\frac{\{1+R_0(t)\}M_{j,u}(t)+R_1(t)B_{j,u}(t)}
         {\{1-R_0(t)\}B_{j,u}(t)}
 \le\frac{C_A}{t}+2R_1(t)=:g(t),
\end{equation}
and $g(t)\to0$. Moreover
$(1-R_0)B_{j,u}\le\mathcal A_{j,u}\le(1+R_0)B_{j,u}$, while the same
tilted bound gives $-(\log B_{j,u})'(t)\le C/t$. Integrating from a
fixed $t_1$ at which $R_0(t_1)\le1/2$, and using $D\le\underline\gamma^{-1}$
and $c\ge c_-$ to bound $B_{j,u}(t_1)$ away from zero, yields uniform
positive constants with $ct^{-C}\le\mathcal A_{j,u}(t)\le2c_+$, hence
\begin{equation}\label{eq:app-logA}
 |\log\mathcal A_{j,u}(t)|\le C\log t
\end{equation}
uniformly for large $t$.

\emph{Step 3: inversion and (E2).} Let
$T_{j,u}(x)=\log Q_j(e^x,u)$. Once $g(t)<1/(2\Gamma)$ the map
$t\mapsto\bar F_j(e^t\mid u)$ is strictly decreasing and invertible, and
\begin{equation}\label{eq:app-inv}
 x=\frac{T_{j,u}(x)}{\xi_j(u)}-\log\mathcal A_{j,u}\{T_{j,u}(x)\}.
\end{equation}
By \eqref{eq:app-logA}, $T_{j,u}(x)\ge c_\gamma x$ for all sufficiently
large $x$. Writing $Q_j(e^x,u)=e^{\xi_j(u)x}\ell_j(e^x,u)$ and
differentiating \eqref{eq:app-inv},
\[
 \frac{\mathrm d}{\mathrm dx}\log\ell_j(e^x,u)
 =\frac{\xi_j(u)^2a_{j,u}\{T_{j,u}(x)\}}
        {1-\xi_j(u)a_{j,u}\{T_{j,u}(x)\}},
 \qquad
 a_{j,u}=\frac{\mathcal A_{j,u}'}{\mathcal A_{j,u}},
\]
so \eqref{eq:app-g} and $T_{j,u}(x)\ge c_\gamma x$ give
\[
 \sup_{j,u}\left|\frac{\mathrm d}{\mathrm dx}\log\ell_j(e^x,u)\right|
 \le\frac{C_E}{x}+4\Gamma^2\overline R_1(c_\gamma x).
\]
For $q>1$ the ratio appearing in (E2) is an average of this derivative
over $[\log s,\log(qs)]$; taking $s\ge(3\alpha_n)^{-1}$ and then all
suprema proves \eqref{eq:transfer-E2}.

\emph{Step 4: tail coupling and (E3$^\uparrow$).} For $t\ge t_0$ set
$f_t(x)=c(x)e^{-t/\gamma(x)}\{1+r(t,x)\}$. Under the coupling of (P4),
with $d=|u-v|$,
\[
 e^{-\delta_t}f_t(X_u)\le f_t(X_v)\le e^{\delta_t}f_t(X_u),
 \qquad
 \delta_t=(L_H+tL_D)d+\zeta_n,
\]
so taking expectations gives
$|\log\bar F_j(e^t\mid v)-\log\bar F_j(e^t\mid u)|\le\delta_t$. By
\eqref{eq:app-g}, after increasing the onset if necessary,
$-\frac{\mathrm d}{\mathrm dt}\log\bar F_j(e^t\mid u)\ge1/(2\Gamma)$,
while $\mathcal A_{j,u}\le2c_+$ and \eqref{eq:app-inv} give
$T_{j,u}(x)\le\Gamma(x+C_Q)$. Evaluating the survival comparison at
$t=T_{j,u}(x)$ and using the inverse-slope bound for the $v$-law,
\[
 |T_{j,v}(x)-T_{j,u}(x)|
 \le2\Gamma\bigl[\{L_H+L_D\Gamma(x+C_Q)\}d+\zeta_n\bigr].
\]
Taking $x=\log(1/w)$ with $w\in[\tau_n,3\alpha_n]$: the lower endpoint
$x\ge x_\alpha\to\infty$ places both quantiles beyond the common onset,
while $x\le2\log(p_nn)$; for $d\le C_0h_n\le\varepsilon/2$ the displaced
point lies in $I_{\varepsilon/2}$. Taking suprema proves
\eqref{eq:transfer-E3}.
\end{proof}

\begin{remark}[The body is never used]
Step~4 controls only $w\le3\alpha_n$. This is exactly the range that
Lemma~\ref{lem:app-tailfreeze} shows to be relevant, and it is why (P4)
constrains the primitive tail alone. No coupling of central quantiles is
required anywhere.
\end{remark}

\subsection{Verifying deficit-mass doubling}

\begin{proposition}[A one-step ratio bound implies (P3)]
\label{prop:app-doubling}
Suppose that, uniformly in $n,j,u$, either $H_{j,u}(0)\ge h_0>0$, or
\begin{equation}\label{eq:app-margin}
 H_{j,u}(z)=a_{j,u}z^{\kappa_{j,n}}L_{j,u}(1/z)\{1+o(1)\},
 \qquad z\downarrow0,
\end{equation}
with $a_{j,u}$ bounded above and away from zero, $0\le\kappa_{j,n}\le
\kappa_+<\infty$, the $o(1)$ uniform, $L_{j,u}$ bounded above and away
from zero at one common point beyond the uniform onset, and
\begin{equation}\label{eq:app-onestep}
 \sup_{n,j,u}\frac{L_{j,u}(t/2)}{L_{j,u}(t)}\le C_L
\end{equation}
for all sufficiently large $t$. Then (P3) holds.
\end{proposition}

\begin{proof}
In the atomic case $H(z)\ge h_0$ for every $z>0$, so
$H(2z)\le c_+\le(c_+/h_0)H(z)$. Otherwise
\[
 \frac{H_{j,u}(2z)}{H_{j,u}(z)}
 =2^{\kappa_{j,n}}\frac{L_{j,u}\{1/(2z)\}}{L_{j,u}(1/z)}\{1+o(1)\},
\]
uniformly bounded by \eqref{eq:app-onestep} and $\kappa_{j,n}\le\kappa_+$.
Taking a fixed small $z_0$, the expansion and the uniform lower bound on
$a_{j,u}L_{j,u}(1/z_0)$ give $H_{j,u}(z_0)\ge h_0'>0$.
\end{proof}

A pure two-sided power margin is not necessary, and is in fact false for
the models below: their corner margins carry powers of $\log(1/z)$ and
factors $\exp\{b\sqrt{\log(1/z)}\}$, which are slowly varying and
therefore satisfy \eqref{eq:app-onestep}.

\subsection{Verification for A1--A3 and B1}

The active set in all four simulation models is $\{1,2,3,4\}$ and the
fibre maximum is attained when every free active PIT coordinate equals
zero. The following boundary calculation is the only model-specific
ingredient; its dimension is at most four, independently of $p_n$.

\begin{lemma}[Gaussian corner]\label{lem:app-corner}
Let $X\sim N_d(\mu,\Sigma)$ with $1\le d\le d_0$, let $\Omega=\Sigma^{-1}$,
and let the parameters range over a compact family of positive definite
covariance matrices and bounded means with
$\min_i(\Omega\mathbf 1)_i\ge\nu_0>0$. Suppose that on a fixed
neighbourhood $N_0$ of the lower corner $D$ is uniformly $C^2$ with
\[
 D\{\Phi(X)\}=\sum_{i=1}^da_i\Phi(X_i)+O(\|\Phi(X)\|^2),
 \qquad 0<a_-\le a_i\le a_+<\infty,
\]
and that $c_*\sum_iu_i\le D(u)\le C_*\sum_iu_i$ on $N_0$ with
$\inf_{u\notin N_0}D(u)>0$. Write $\kappa=\mathbf 1^\top\Omega\mathbf 1$,
$\nu=\Omega\mathbf 1$, $b=\mathbf 1^\top\Omega\mu$ and
$r_z=-\Phi^{-1}(z)$. Then, uniformly over the family,
\begin{equation}\label{eq:app-corner}
 \Pp[D\{\Phi(X)\}\le z]
 =C(\mu,\Sigma,a)\,z^{\kappa}\{\sqrt{2\pi}r_z\}^{\kappa-d}e^{-br_z}
 \{1+O(r_z^{-1})\},
\end{equation}
with
$C(\mu,\Sigma,a)=|\Sigma|^{-1/2}e^{-\mu^\top\Omega\mu/2}
\prod_i\Gamma(\nu_i)a_i^{-\nu_i}/\Gamma(\kappa+1)$.
The conclusion persists after multiplication by a positive weight that
is uniformly Lipschitz at the corner and bounded there above and away
from zero. In particular (P3) holds.
\end{lemma}

\begin{proof}
The density of $U=\Phi(X)$ is
\[
 f_U(u)=|\Sigma|^{-1/2}
 \exp\bigl\{-\tfrac12z(u)^\top(\Omega-I)z(u)+z(u)^\top\Omega\mu
 -\tfrac12\mu^\top\Omega\mu\bigr\},
 \quad z_i(u)=\Phi^{-1}(u_i).
\]
Set $u_i=\delta x_i$ and $r=r_\delta$. Uniformly on compact subsets of
$(0,\infty)^d$,
$\Phi^{-1}(\delta x_i)=-r+r^{-1}\log x_i+o(r^{-1})$ and
$e^{-r^2/2}=\sqrt{2\pi}r\delta\{1+O(r^{-2})\}$, so with $\nu=\Omega\mathbf 1$
\[
 f_U(\delta x)\delta^d=|\Sigma|^{-1/2}e^{-\mu^\top\Omega\mu/2}
 \delta^{\kappa}\{\sqrt{2\pi}r\}^{\kappa-d}e^{-br}
 \prod_ix_i^{\nu_i-1}\{1+o(1)\}.
\]
The two-sided bound on $D$ confines $\{D(\delta x)\le\delta\}$ to a
fixed rescaled simplex and excludes points outside $N_0$. Choose
$0<\theta<\tfrac12\min\{1,\nu_0,\underline\lambda\}$ with
$\underline\lambda=\inf\lambda_{\min}(\Omega)$, and split the simplex into
$\mathcal B_r=\{x:\min_ix_i<e^{-\sqrt r}\}$ and its complement. On
$\mathcal B_r$, Mills bounds and convexity of the Gaussian quadratic form
give the envelope $C\prod_ix_i^{\theta-1}$, and a one-dimensional beta
integral bounds the normalized mass of $\mathcal B_r$ by
$Ce^{-\theta\sqrt r}=o(r^{-1})$. On $\mathcal B_r^c$ we have
$|\log x_i|\le\sqrt r$ from below and a uniform bound from above, and the
uniform normal-quantile expansion
\[
 \Phi^{-1}(\delta x_i)=-r+\frac{\log x_i}{r}
 +O\!\left(\frac{1+|\log x_i|+(\log x_i)^2+\log r}{r^3}\right)
\]
shows that the quadratic Gaussian term contributes a relative error
$r^{-1}$ times a fixed polynomial in $1+\sum_i|\log x_i|$, whose
Dirichlet log moments are finite uniformly because $\nu_i\ge\nu_0$. The
quadratic remainder in $D$ moves the limiting simplex boundary by
$O(\delta)=o(r^{-1})$. Integrating the density expansion and using
\[
 \int_{\{\sum_ia_ix_i\le1\}}\prod_ix_i^{\nu_i-1}\,\mathrm dx
 =\frac{\prod_i\Gamma(\nu_i)a_i^{-\nu_i}}{\Gamma(\kappa+1)}
\]
gives \eqref{eq:app-corner}. The factor multiplying $z^\kappa$ is
slowly varying and obeys \eqref{eq:app-onestep}, so
Proposition~\ref{prop:app-doubling} applies. A positive weight
$w(u)=w(0)+O(\|u\|_1)$ changes the rescaled integrand by a relative
$O(\delta)$ and therefore only multiplies the leading constant.
\end{proof}

For the AR(1) covariance of the four active Gaussian coordinates at
$\rho=1/4$, the corner exponents are
\begin{equation}\label{eq:app-kappas}
 \kappa_j=
 \begin{cases}
  34/15, & j=1,4,\\
  41/15, & j=2,3,\\
  14/5+\lambda_j^2/(1-\lambda_j^2), & j\notin\{1,2,3,4\},
 \end{cases}
 \qquad \lambda_j=\rho^{\,j-4}\ (j\ge5),
\end{equation}
so $34/15\le\kappa_j\le43/15$ uniformly; in B1 a proxy coordinate is
independent of the active block and has $\kappa_j=14/5$. All components
of $\Omega\mathbf 1$ are bounded away from zero. The reciprocal deficit
has the required expansion: in A1 and A2 this follows by expanding the
exponential index surface at the active corner, and in A3 the relevant
first derivatives of its exponent are $0.8+0.20u$, $0.8+0.15u$ or $0.8$,
hence uniformly at least $0.8$. In the A family the weight $c$ is
identically one; in B1 the bounded latent scale factor, integrated over
the conditional Gaussian law of $F$, has a uniform expansion
$w_{j,u}(z)=w_{j,u}(0)+O(\|z\|_1)$ with $0<w_-\le w_{j,u}(0)\le w_+$, so
the weighted clause of Lemma~\ref{lem:app-corner} applies.

The remaining primitive assumptions hold as follows. The conditional
generators are strictly increasing and atomless, which gives (P1). For
A1--A3, $q(s,u)=s^{\gamma(u)}/[1+\exp\{v(u)-s\}]$, and both the survival
remainder and its log-level derivative decay faster than every
exponential in $t=\log y$; for B1,
$q(s,u,F)=s^{\gamma(u)}\exp\{A(F)-1/(2s)\}$ with
$A(F)=\{\Phi(F)-1/2\}/2$, and inversion gives
$\Pp(Y>y\mid U=x)=c_B(x)y^{-1/\gamma(x)}\{1+O(y^{-2})\}$ with
$c_B(x)=\E[e^{A(F)/\gamma}\mid U=x]$, the same order holding for the
log-level derivative. This is (P2). A common-residual coupling of the
conditional Gaussian laws moves the other PIT coordinates by
$O(|u-v|)$ on $I_{\varepsilon/2}$, which gives the $1/\gamma$ clause of
(P4); the $c$ and $\log(1+r)$ clauses follow by differentiating the
bounded posterior exponential moments, with $\zeta_n=0$. The projected
profiles are exactly
\[
 \xi_j(u)=
 \begin{cases}
  \tfrac12e^{-u}, & \text{A1, B1},\ j\le4,\\
  \tfrac12e^{-u/2}, & \text{A2},\ j\le4,\\
  \tfrac12e^{-0.8u}, & \text{A3},\ j\le4,\\
  \tfrac12, & j>4.
 \end{cases}
\]

\begin{corollary}[The simulation models satisfy the estimation
conditions]\label{cor:app-models}
For A1--A3 and B1, conditions (E1), (E2) and (E3$^\uparrow$) hold along
every sequence
\[
 \alpha_n=n^{-a},\qquad h_n=\tfrac12n^{-b},\qquad p_n\le n^{C},
\]
with fixed $0<a<1$ and $b,C>0$, with deterministic bias orders
\[
 \omega_\xi(1/n)=O(n^{-1}),
 \qquad
 b_n(3\alpha_n)=O(1/\log n),
 \qquad
 \omega_n^{\uparrow}(C_0h_n,\tau_n,\alpha_n)=O(n^{-b}\log n).
\]
If in addition $b<1/2$ and $a+b<1$, then
$\log(p_nn)=o(nh_n^2)$ and $\log(p_nn)=o(n\alpha_nh_n)$, so the
population gaps of A1--A3 and B1 being fixed and positive,
Theorems~\ref{thm:score} and~\ref{thm:sis} apply. This covers all nine
exponent pairs used by the aggregated screen: intersecting the nine
separation events, every active coordinate has rank at most $s$ and
every inactive coordinate rank at least $s+1$ at every setting, so
taking the minimum of the nine ranks preserves the same strict
separation, and a finite union bound makes the intersection probability
tend to one.
\end{corollary}

\begin{remark}[Finite-design honesty]\label{rem:app-finite}
Corollary~\ref{cor:app-models} is asymptotic. It does not certify that
the $o(\Delta_{\min})$ biases are numerically small at the design used
in Section~\ref{sec:simulation}: there $\log(1/\alpha)\simeq2.28$, the
lower endpoint entering (E2) is only
$x_\alpha=\log\{(3\alpha)^{-1}\}\simeq1.18$, and $h\log(pn)$ is much
larger than the population gaps under the available uniform upper
bounds. The finite-sample tables are empirical evidence about that
regime; they are not a numerical verification of (E2), (E3$^\uparrow$)
or of the concentration bound.
\end{remark}
